\documentclass[utf8]{article} %中文文档类
\usepackage[left=2.50cm, right=2.50cm, top=2.50cm, bottom=2.50cm]{geometry} % 页边距
\usepackage{indentfirst} % 首行缩进
\usepackage{fancyhdr,fancybox} % 设置页眉、页脚
\usepackage{ctexcap} % 标题是中文的
\usepackage{helvet} % 用来指定beamer使用的字体
\usepackage{hyperref} % bookmarks 
\usepackage{multicol} % 分栏
\usepackage{cite} % 引用

\usepackage{graphicx} % 图片
\usepackage{subcaption} % 子图
\usepackage{multirow} % 表格
\usepackage{float} % 浮动体
\usepackage{booktabs} % 三线表
\usepackage{longtable} % 长表格
\usepackage{diagbox} % 表头斜线
\usepackage{multirow} % 表格合并
\usepackage{makecell} % 表格换行
\usepackage{enumitem} % 枚举环境宏包
\usepackage{framed} % 边框包

\usepackage{amsmath, amsfonts, amssymb} % 数学公式、符号
\numberwithin{equation}{section}   % 编号从一级标题开始
\usepackage[english]{babel} % 数学公式标准
\usepackage{bm} % 加粗方程字体 
\usepackage{caption} % 标题
\usepackage{color} % 配色
\usepackage[braket]{qcircuit} % 启用狄拉克符号支持, 用于绘制电路图
\usepackage{physics} % 量子力学中的符号加载

\newtheorem{Def}{\hskip 2em {\color{green} 定义}}[subsection]
\newtheorem{Thm}{\hskip 2em {\color{blue} 定理}}[subsection]
\newtheorem{Lem}{\hskip 2em {\color{magenta} 引理}}[subsection]
\newtheorem{Prop}{\hskip 2em {\color{cyan} 命题}}[subsection]
\newtheorem{Note}{\hskip 2em {\color{red} 注}}[subsection]
\newtheorem{Cor}{\hskip 2em {\color{cyan} 推论}}[subsection]
\newtheorem{Eg}{\hskip 2em {\color{black} 例}}[subsection]
\newenvironment{Proof}{{\noindent\it Proof:}\quad}{\hfill $\sharp $\par}
\allowdisplaybreaks

% 页眉页脚设置
\usepackage{fancyhdr}
\usepackage{lastpage} % 获得总页数
\pagestyle{fancy}
\fancyfoot{}
\lhead{\textit{Modern Framework For Quantum Algorithms}}\chead{}\rhead{\textit{\thepage\ of \pageref{LastPage}}}
\lfoot{}\cfoot{}\rfoot{}
\renewcommand{\headrulewidth}{0.4pt}

% 标题设置
\title{\textbf{Module3 现代量子算法统一框架}}
\date{\today}
\author{\textbf{STU\quad 牛晓铭}}

\hypersetup{
	colorlinks=true,
	linkcolor=black
} % 使目录为黑色字
\begin{document}  
\maketitle
\renewcommand{\contentsname}{} 
\tableofcontents
\thispagestyle{empty}
	
\newpage
\setcounter{page}{1}

\section{引言: 从原语到统一框架}

Module2 中的四类\textbf{原语(primitive)}有一个共同的困难: 它们都要处理\textbf{非酉的对象(non-unitary)}. Trotter 型\textbf{哈密顿量模拟(Hamiltonian Simulation)}
对时间依赖与病态哈密顿量的误差控制较差(步数 $L = O(t^{1+1/p}\varepsilon^{-1/p})$); \textbf{量子相位估计(QPE)} 需要访问
$U^{2^j}$, 而一般的\textbf{矩阵函数(matrix function)}(如 $A^{-1}$、$e^{-iAt}$)并不对应直接的\textbf{酉预言机(unitary oracle)}; \textbf{振幅放大(amplitude amplification)}需要
构造\textbf{反射算子(reflection operator)}, 其实现依赖具体问题. 本模块介绍\textbf{现代统一框架}. 开头先简练复习 Module2 的非感知振幅放大 (OAA),
随后沿"输入矩阵 → 变换矩阵 → 逼近函数"的线索层层递进:

\begin{enumerate}
    \item \textbf{LCU(线性组合酉)}: 用 prepare/select oracle 实现矩阵的线性组合,
    是最简单的"输入矩阵"方式;
    \item \textbf{Block-encoding(块编码)}: 把任意(非酉)矩阵 $A$ 嵌入一个更大的\textbf{酉算子(unitary operator)}
    $U_A$ 的左上角, 从而在量子计算机上"输入"任意矩阵;
    \item \textbf{Qubitization 与 Chebyshev 多项式}: 把 block-encoding 中每个本征(奇异)方向上的
    "反射"转化为"旋转", 一步实现\textbf{Chebyshev 多项式(Chebyshev polynomials)} $T_k(A)$;
    后者同时是\textbf{多项式逼近(polynomial approximation)}的基本工具箱, 连接 qubitization 与 QSP;
    \item \textbf{QSP/QET 与 QSVT}: 用相位因子序列把任意(满足奇偶性条件的)多项式
    $P(A)$ 直接"写入"量子电路 —— 这是现代量子算法的\textbf{统一实现层}.
\end{enumerate}

\section{OAA 复习: 与输入态无关的振幅放大}

在展开 LCU、block-encoding 等构件之前, 先简练复习 Module2 中学过的
\textbf{非感知振幅放大(oblivious amplitude amplification, OAA)}. 它处理的对象是
"编码—后选择"两步: 若某个酉算子 $U$ 被编码为
\begin{equation}
    V\ket{0^a}\ket{\psi} = \frac{1}{\alpha}\ket{0^a}\, U\ket{\psi} + \ket{\Phi^\perp},
    \qquad \big(\bra{0^a}\otimes I_n\big)\ket{\Phi^\perp} = 0,
\end{equation}
则单次后选择成功概率 $p = 1/\alpha^2$ 与输入态 $\ket{\psi}$ 无关.

\begin{Def}[OAA 走算子]
    记成功子空间投影 $\Pi = \ket{0^a}\bra{0^a}\otimes I_n$ 与反射
    $Z_\Pi = I - 2\Pi$. \textbf{OAA 走算子(walk operator)}定义为
    \begin{equation}
        Q_{\mathrm{obl}} := -\, V\, Z_\Pi\, V^\dagger\, Z_\Pi.
    \end{equation}
    记 $\sin\theta = 1/\alpha$, 则 $Q_{\mathrm{obl}}$ 在二维不变子空间
    $\mathrm{span}\{\ket{0^a}U\ket{\psi}, \ket{\Phi^\perp}\}$ 上是角度
    $2\theta$ 的旋转, 与输入态 $\ket{\psi}$ 无关.
\end{Def}

\begin{Note}[几何与关键性质]
    经 $k$ 轮放大后, 成功子空间上的振幅为
    \begin{equation}
        \big(\Pi \otimes I_n\big)\, Q_{\mathrm{obl}}^k\, V\ket{\psi}
        = U\ket{\psi}\cdot \sin\big((2k+1)\theta\big),
    \end{equation}
    即振幅(而非概率)随 $k$ 线性增长, 经 $O(\theta^{-1}) = O(\alpha)$ 轮逼近 $1$.
    三个关键性质: ① 旋转角与输入态无关 —— 这是与普通振幅放大(需要制备并反制备
    当前输入态)的根本区别, 使 OAA 可以安全嵌套进多步算法; ② $\alpha = 2$ 时一轮即可
    把成功概率从 $1/4$ 提高到 $1$; ③ 对目标块非酉的一般矩阵, 不能使用 oblivious
    版本, 应回到 Module2 依赖输入态的标准振幅放大.
\end{Note}

\begin{Note}[与后面的联系]
    上面的编码 $V\ket{0^a}\ket{\psi}$ 正是下文 LCU 引理与 block-encoding 的雏形;
    OAA 走算子 $Q_{\mathrm{obl}}$ 与后面 qubitization 的迭代算子 $U_A Z U_A^\dagger Z$
    具有相同的"编码—反射—解编码—反射"结构 —— 本模块将在这条线索上系统化.
\end{Note}

\section{LCU: 线性组合的量子实现}

\textbf{线性组合酉(linear combination of unitaries, LCU)}是量子算法中实现一般线性算子的一种
基本方法: 把一个通常非酉的算子 $A$ 表示为若干易实现的酉算子的线性组合, 再借助辅助寄存器、
受控酉操作与量子干涉, 把 $A$ 的作用嵌入到一个更大的酉演化中.

\subsection{基本设定}

设目标算子 $A$ 可以写成
\begin{equation}
    A = \sum_{j=0}^{L-1} \alpha_j U_j,
    \qquad U_j^\dagger U_j = I,
\end{equation}
即每个 $U_j$ 都是酉算子. 先假设系数 $\alpha_j \ge 0$, 记归一化常数
\begin{equation}
    \alpha := \sum_{j=0}^{L-1} \alpha_j.
\end{equation}
若原始系数是复数 $c_j = |c_j| e^{i\phi_j}$, 则把相位吸收进对应的酉算子
$\widetilde U_j := e^{i\phi_j} U_j$, 有 $c_j U_j = |c_j|\,\widetilde U_j$ ——
因此总可以把 LCU 写成系数非负的形式. LCU 的目标是实现
\begin{equation}
    A\ket{\psi} = \sum_j \alpha_j U_j \ket{\psi};
\end{equation}
由于 $A$ 一般不是酉算子, 不能直接作为量子门实现, 需要嵌入到一个更大的酉算子中.

\subsection{PREPARE 与 SELECT 算子}

引入辅助寄存器, 定义酉算子 $O_{prep}$(\textbf{prepare oracle})满足
\begin{equation}
    O_{prep}\ket{0^a} = \sum_{j=0}^{L-1} \sqrt{\frac{\alpha_j}{\alpha}}\,\ket{j}
    =: \ket{\chi},
    \qquad
    \sum_j \frac{\alpha_j}{\alpha} = 1
    \;\Longrightarrow\; \langle\chi|\chi\rangle = 1,
\end{equation}
系数以\textbf{平方根}编码进振幅, 是因为测量概率等于振幅的模方. 受控
\textbf{select oracle} 定义为
\begin{equation}
    O_{sel} := \sum_{j=0}^{L-1} \ket{j}\bra{j} \otimes U_j,
    \qquad
    O_{sel}\ket{j}\ket{\psi} = \ket{j} U_j \ket{\psi},
\end{equation}
即辅助态 $\ket{j}$ 决定数据寄存器上执行哪个 $U_j$.

\subsection{LCU 引理}

定义整个 LCU 酉算子
\begin{equation}
    W := (O_{prep}^\dagger \otimes I_n)\, O_{sel}\, (O_{prep} \otimes I_n),
\end{equation}
它对应电路序列 $O_{prep} \to O_{sel} \to O_{prep}^\dagger$.

\begin{Thm}[LCU 引理] \label{thm:lcu}
    对任意 $\ket{\psi}$,
    \begin{equation}
        \big(\bra{0^a}\otimes I_n\big) W \big(\ket{0^a}\otimes I_n\big) = \frac{A}{\alpha},
    \end{equation}
    \begin{equation}
        W \ket{0^a}\ket{\psi} = \frac{1}{\alpha}\ket{0^a} A \ket{\psi} + \ket{\Phi_\psi^\perp},
    \end{equation}
    其中 $\big(\bra{0^a}\otimes I_n\big)\ket{\Phi_\psi^\perp} = 0$. 第一项为
    \textbf{成功分量}, 第二项为\textbf{失败分量}; 测量辅助寄存器得到 $\ket{0^a}$ 时,
    数据寄存器态为 $A\ket{\psi}/\norm{A\ket{\psi}}$ —— LCU 是\textbf{概率性}实现一般
    线性算子 $A$ 的方法.
\end{Thm}

\begin{Proof}
    由 $\bra{0^a}O_{prep}^\dagger = \sum_j \sqrt{\alpha_j/\alpha}\,\bra{j}$ 与 SELECT 的正交性
    $(\bra{j}\otimes I)O_{sel}(\ket{k}\otimes I) = \delta_{jk} U_k$:
    \begin{align}
        \big(\bra{0^a}\otimes I\big) W \big(\ket{0^a}\otimes I\big) &=\big(\bra{0^a}O_{prep}^\dagger\otimes I\big)O_{sel}\big(O_{prep}\ket{0^a} \otimes I\big)\\
        &= \sum_{j,k} \sqrt{\frac{\alpha_j}{\alpha}}\sqrt{\frac{\alpha_k}{\alpha}}
           \big(\bra{j}\otimes I\big)O_{sel}\big(\ket{k}\otimes I\big)
        \nonumber\\
        &= \sum_j \frac{\alpha_j}{\alpha} U_j
         = \frac{1}{\alpha}\sum_j \alpha_j U_j
         = \frac{A}{\alpha}.
    \end{align}
    辅助比特数只依赖 $\log L$, 与线性组合的项数成对数关系 —— 这是 LCU 高效性的关键.
\end{Proof}

\begin{Note}[成功概率与振幅放大]
    成功分量的振幅为 $A\ket{\psi}/\alpha$, 故成功概率为
    \begin{equation}
        p_{\mathrm{success}} = \frac{\norm{A\ket{\psi}}^2}{\alpha^2}.
    \end{equation}
    直接重复测量平均需 $(\alpha/\norm{A\ket{\psi}})^2$ 次; 结合振幅放大(Module2)降为
    $O(\alpha/\norm{A\ket{\psi}})$. 特别地, 若 $A$ 本身是酉算子, 则
    $\norm{A\ket{\psi}} = 1$, 成功概率 $= 1/\alpha^2$ 与输入态无关, 可直接用第 2 节的
    \textbf{OAA} 放大, 重复次数从 $O(\alpha^2)$ 降为 $O(\alpha)$.
    因此 \textbf{LCU 构造所需的线性组合, OAA 放大其成功振幅} —— 两者经常组合使用.
\end{Note}

\subsection{LCU 的数学本质}

LCU 可以概括为三步:
\begin{equation}
    \text{PREPARE} \;\to\; \text{SELECT} \;\to\; \text{PREPARE}^\dagger,
    \qquad
    \text{superposition} \;\to\; \text{controlled unitaries} \;\to\; \text{interference}.
\end{equation}
PREPARE 把线性组合的系数编码到辅助寄存器振幅中, SELECT 在不同辅助分支上执行不同的酉算子,
PREPARE$^\dagger$ 使不同量子路径重新发生干涉, 在成功分支中产生所需的线性组合.

\subsection{LCU 与 Hamiltonian 模拟}

若哈密顿量可以写成 $H = \sum_j \alpha_j U_j$, 则可借 LCU 表示 $H$. Hamiltonian 模拟的目标是
实现 $e^{-iHt}$, 由 \textbf{Taylor 展开}
\begin{equation}
    e^{-iHt} = \sum_{k \ge 0} \frac{(-it)^k}{k!} H^k,
    \qquad
    H^k = \sum_{j_1,\dots,j_k} \alpha_{j_1}\cdots\alpha_{j_k}\, U_{j_1}\cdots U_{j_k},
\end{equation}
而酉算子之积 $U_{j_1}\cdots U_{j_k}$ 仍是酉算子, 故对 Taylor 级数有限截断后
\begin{equation}
    e^{-iHt} \approx \sum_{\ell} \beta_\ell V_\ell,
    \qquad V_\ell \text{ 均酉},
\end{equation}
Hamiltonian 模拟转化为另一个 LCU 问题 —— 这是\textbf{截断 Taylor 级数 + LCU} 型
Hamiltonian 模拟算法的基础.

\subsection{与 block-encoding 的联系}

LCU 的核心结果 $\big(\bra{0^a}\otimes I\big) W \big(\ket{0^a}\otimes I\big) = A/\alpha$
意味着从矩阵结构上看
\begin{equation}
    W = \begin{pmatrix} A/\alpha & \ast \\ \ast & \ast \end{pmatrix},
\end{equation}
即一般的非酉算子 $A$ 被编码进一个更大的酉矩阵 $W$ 的一个块中 —— 这正是下一节
\textbf{block-encoding} 的基本思想. 因此 \textbf{LCU 是构造 block-encoding 的一种重要方法}.

\subsection{两个例子}

\begin{Eg}[两项 LCU: 两个酉之和]
    对 $A = aU_0 + bU_1$($a,b > 0$), $\alpha = a + b$, prepare oracle 为
    $O_{prep}\ket{0} = \sqrt{a/\alpha}\,\ket{0} + \sqrt{b/\alpha}\,\ket{1}$
    (特例 $a = b$ 时 $O_{prep} = H$), select oracle 为受控 $U_0$, $U_1$: 电路为
    \begin{equation*}
    \Qcircuit @C=1em @R=.8em {
        \lstick{\ket{0}} & \gate{O_{prep}} & \ctrl{1} & \ctrl{1} & \gate{O_{prep}^\dagger} & \qw & \rstick{\text{测量}} \qw \\
        \lstick{\ket{\psi}} & \qw & \gate{U_0} & \gate{U_1} & \qw & \qw & \qw
    }
    \end{equation*}
    经 SELECT 后为 $\sqrt{a/\alpha}\,\ket{0} U_0\ket{\psi} + \sqrt{b/\alpha}\,\ket{1} U_1\ket{\psi}$,
    再作用 $O_{prep}^\dagger$, 辅助态回到 $\ket{0}$ 的分量即
    $\big(a U_0 + b U_1\big)\ket{\psi}/\alpha = A\ket{\psi}/\alpha$ ——
    两条量子路径在 $O_{prep}^\dagger$ 处发生干涉, 产生所需的线性组合.
    对两个(下一节意义的)block-encoded 矩阵 $A + B$, 只需把受控门换成受控的 $U_A$, $U_B$
    即可(线性组合的非酉推广).
\end{Eg}

\begin{Eg}[Chebyshev 线性组合: \textbf{矩阵多项式(matrix polynomial)}]
    后面 qubitization 与 Chebyshev 一节将给出 $T_k(A)$ 的 block-encoding
    ($U_k := (U_A Z)^k$); 届时与 LCU 结合,
    对 $f(A) = \sum_k \alpha_k T_k(A)$ 得到 $f(A)$ 的 $(\alpha, 2a)$-block-encoding
    ($\alpha = \sum_k |\alpha_k|$, 记号见下一节 block-encoding).
    这是 HHL 与基态滤波等问题的经典实现路线.
\end{Eg}

\subsection{LCU 的局限}

\begin{Note}[LCU 的局限]
    LCU 需要实现 select 与 prepare oracle, 其中 prepare 的振幅 $\sqrt{\alpha_j/\alpha}$
    需要多比特受控旋转; 直接实现时代价随 $L$ 线性增长. 此外 $f(A)$ 的
    1-范数 $\alpha$ 通常较大(如 Chebyshev 展开的 $\alpha$ 随度数增长).
    后面的 QSP/QET 正是为了消除这两个困难而提出.
\end{Note}

\section{Block-encoding: 非酉矩阵的输入模型}

\textbf{块编码(block-encoding)}是现代量子算法中表示一般线性算子的基本框架之一. 量子电路本身
必须由酉算子构成, 而很多实际问题中的目标矩阵 $A$ 并不是酉矩阵, 不能直接作为量子门实现.
Block-encoding 通过引入辅助量子比特, 将经过适当归一化的矩阵 $A/\alpha$ 嵌入到一个更高维
酉矩阵的某个子块中:
\begin{equation}
    A \;\longrightarrow\; U_A = \begin{pmatrix} A/\alpha & \ast \\ \ast & \ast \end{pmatrix},
    \qquad U_A^\dagger U_A = I.
\end{equation}
它建立了 \textbf{general matrix} $\longrightarrow$ \textbf{unitary quantum circuit} 之间的桥梁.

\subsection{为什么需要 block-encoding}

设 $A \in \mathbb{C}^{N \times N}$ 是希望在量子计算机上处理的矩阵. 一般情况下
$A^\dagger A \ne I$, 即 $A$ 不是酉矩阵; 而封闭量子系统中的确定性量子演化必须由酉算子描述
($U^\dagger U = I$), 所以一般不能直接执行 $\ket{\psi} \to A\ket{\psi}$.
Block-encoding 的基本策略不是直接把 $A$ 作为量子门, 而是寻找一个更大的酉算子 $U_A$,
使得 $A$ 出现在 $U_A$ 的某个矩阵块中.

\subsection{输入模型与矩阵查询 oracle}

在量子计算机上处理矩阵 $A \in \mathbb{C}^{N \times N}$($N = 2^n$), 首先要回答的问题是
\textbf{输入模型(input model)}: 如何访问 $A$ 的信息? 一种方式是 $e^{iA}$(若 $A$ \textbf{非厄米(non-Hermitian)}, 可用\textbf{膨胀(dilation)}方法),
其优点是可以直接用 Trotter 等简单电路实现; 但更一般的模型是本节介绍的 block-encoding.
若 $A$ 是\textbf{稠密矩阵(dense matrix)}, 任何输入模型都需要指数多资源, 因此通常假设 $A$ 是\textbf{$s$-稀疏(s-sparse)}的
(每行每列至多 $s$ 个非零元), 且非零元的位置与数值可以高效查询.

\begin{Def}[矩阵查询 oracle]
    记 $\norm{A}_{\max} := \max_{ij}|A_{ij}|$. \textbf{矩阵查询预言机(matrix query oracle)} 形如
    \begin{equation}
        O_A \ket{0}\ket{i}\ket{j}
        = A_{ij}\ket{0}\ket{i}\ket{j} + \sqrt{1 - |A_{ij}|^2}\,\ket{1}\ket{i}\ket{j},
    \end{equation}
    即以 $A_{ij}$ 为振幅对信号比特做受控旋转(若 $\norm{A}_{\max} \ge 1$,
    先考虑缩放后的矩阵 $A/\alpha$). 更初级的模型是先把 $A_{ij}$ 算进工作寄存器
    (定点表示), 再通过受控旋转与 uncomputation 转化为上式.
\end{Def}

\subsection{Block-encoding 的定义}

设系统寄存器由 $n$ 个量子比特组成, 辅助寄存器由 $a$ 个量子比特组成,
$\ket{0^a} = \ket{0}^{\otimes a}$.

\begin{Def}[Block-encoding] \label{def:be}
    给定 $n$ 比特矩阵 $A$, 若存在 $\alpha > 0$, 非负整数 $a$ 与 $(a+n)$ 比特酉算子 $U_A$ 使得
    \begin{equation}
        \norm{A - \alpha \big(\bra{0^a}\otimes I\big) U_A \big(\ket{0^a}\otimes I\big)} \le \varepsilon,
    \end{equation}
    则称 $U_A$ 是 $A$ 的 $(\alpha, a, \varepsilon)$-\textbf{block-encoding};
    当 $\varepsilon = 0$ 时称 \textbf{exact block-encoding}, 简记 $(\alpha, a)$-block-encoding, 此时
    \begin{equation}
        \big(\bra{0^a}\otimes I\big) U_A \big(\ket{0^a}\otimes I\big) = A/\alpha.
    \end{equation}
    记 $\mathrm{BE}_{\alpha,a}(A, \varepsilon)$ 为所有这样的 $U_A$ 的集合.
    直观地说, $U_A$ 的左上 $N \times N$ 块(由 $\ket{0^a}$ 标志)正是 $A/\alpha$.
    三个参数分别刻画\textbf{归一化因子(normalization factor)} $\alpha$、\textbf{辅助比特数} $a$、
    与\textbf{编码误差} $\varepsilon$.
\end{Def}

\begin{Note}[投影形式与块矩阵形式]
    定义\textbf{成功子空间投影} $\Pi = \ket{0^a}\bra{0^a}\otimes I$. 定义可等价写成
    \begin{equation}
        \Pi U_A \Pi = \ket{0^a}\bra{0^a} \otimes \frac{A}{\alpha},
        \qquad
        U_A = \begin{pmatrix} U_{00} & U_{01} \\ U_{10} & U_{11} \end{pmatrix},
        \quad U_{00} = \big(\bra{0^a}\otimes I\big) U_A \big(\ket{0^a}\otimes I\big) = \frac{A}{\alpha}.
    \end{equation}
    即在辅助寄存器处于 $\ket{0^a}$ 的子空间中, $U_A$ 的有效作用等于 $A/\alpha$ ——
    这是 "block encoding" 这一名称的来源.
\end{Note}

\begin{Prop}[测量与成功概率]
    设 $U_A$ 是 $A$ 的 $(\alpha, a)$-block-encoding, 对任意 $n$ 比特态 $\ket{b}$,
    \begin{equation}
        U_A \ket{0^a}\ket{b} = \frac{1}{\alpha}\ket{0^a} A\ket{b} + \ket{\Phi^\perp},
        \qquad \big(\bra{0^a}\otimes I\big)\ket{\Phi^\perp} = 0.
    \end{equation}
    第一项位于辅助寄存器为 $\ket{0^a}$ 的\textbf{成功子空间}, 第二项位于其正交补. 测量
    $a$ 个辅助比特, 全部得到 $0$ 的概率为
    \begin{equation}
        p(0^a) = \frac{1}{\alpha^2}\norm{A\ket{b}}^2,
    \end{equation}
    且测量后系统寄存器处于归一化的 $A\ket{b}/\norm{A\ket{b}}$. 若 $\alpha$ 过大,
    成功概率会指数小 —— 因此 $\alpha$ 是 block-encoding 的重要代价参数.
\end{Prop}

\begin{Note}[归一化因子 $\alpha$ 的作用]
    由于 $U_A$ 是酉矩阵, 其任意子块的算子范数不能超过 $1$, 因此必须满足
    \begin{equation}
        \norm{\frac{A}{\alpha}} \le 1
        \quad\Longrightarrow\quad
        \alpha \ge \norm{A}
        \quad(\text{通常取谱范数 } \norm{A} = \sigma_{\max}(A)).
    \end{equation}
    因此 $\alpha$ 是把 $A$ 缩放到单位范数以内的归一化常数; 算法设计中希望 $\alpha$ 尽可能
    接近 $\norm{A}$, 因为过大的 $\alpha$ 会降低成功概率、增加后续算法复杂度.
\end{Note}

\subsection{若干构造}

\begin{Eg}[对角矩阵]
    对角矩阵 $A = \mathrm{diag}(A_{00}, \dots, A_{N-1,N-1})$ 的查询 oracle 简化为
    $O_A\ket{0}\ket{i} = A_{ii}\ket{0}\ket{i} + \sqrt{1-|A_{ii}|^2}\ket{1}\ket{i}$,
    直接验证 $\bra{0}\bra{i} O_A \ket{0}\ket{j} = A_{ii}\delta_{ij}$,
    故 $O_A$ 是 $(1,1)$-block-encoding.
\end{Eg}

\begin{Thm}[\textbf{$s$-稀疏矩阵(s-sparse matrix)}的 block-encoding]
    对 $s = 2^{s'}$ 稀疏的矩阵, 假设可以访问行索引 oracle
    $O_c\ket{\ell}\ket{j} = \ket{\ell}\ket{c(j,\ell)}$(第 $j$ 列第 $\ell$ 个非零元所在行)与
    矩阵元 oracle $O_A$(见 4.1 节), 令 $D = I_{n-s'} \otimes H^{\otimes s'}$
    为扩散算子($D\ket{0^{s'}} = \frac{1}{\sqrt{s}}\sum_{\ell}\ket{\ell}$), 则电路
    \begin{equation*}
    \Qcircuit @C=.9em @R=.8em {
        \lstick{\ket{0}} & \qw & \qw & \qw & \qw & \qw & \rstick{\text{测量}} \qw \\
        \lstick{\ket{0^{s'}}} & \gate{D} & \gate{O_A} & \gate{D} & \qw & \qw & \qw \\
        \lstick{\ket{j}} & \qw & \qw & \gate{O_c} & \qw & \qw & \rstick{\ket{i}} \qw
    }
    \end{equation*}
    定义 $U_A$, 则 $U_A \in \mathrm{BE}_{s,\, s'+1}(A)$, 即
    \[ \big(\bra{0}\bra{0^{s'}}\otimes I_n\big) U_A \big(\ket{0}\ket{0^{s'}}\otimes I_n\big) = A/s. \]
\end{Thm}

\begin{Proof}
    对源态依次施加 $D$, $O_A$, $O_c$:
    \[ \ket{0}\ket{0^{s'}}\ket{j} \;\xrightarrow{D}\;
       \frac{1}{\sqrt{s}}\sum_{\ell} \ket{0}\ket{\ell}\ket{j}
       \;\xrightarrow{O_A,O_c}\;
       \frac{1}{\sqrt{s}}\sum_{\ell} \Big(A_{c(j,\ell),j}\ket{0}
       + \sqrt{1-|A_{c(j,\ell),j}|^2}\ket{1}\Big)\ket{\ell}\ket{c(j,\ell)}. \]
    对目标态 $\bra{0}\bra{0^{s'}}\bra{i}$ 同样施加 $D$ 后与源态求内积
    (注意 $D^\dagger = D$), 非零贡献仅来自 $\ell$ 使 $c(j,\ell) = i$ 的项,
    共 $\frac{1}{s} A_{ij}$ —— 即左上块为 $A/s$, 证毕.
\end{Proof}

\begin{Eg}[SVD 构造: $(1,1)$-block-encoding 的一般性]
    对任意满足 $\norm{A} \le 1$ 的矩阵, 由\textbf{奇异值分解(SVD)} $A = W\Sigma V^\dagger$ 可构造
    \begin{equation}
        U_A = \begin{pmatrix} W & 0 \\ 0 & I_n \end{pmatrix}
        \begin{pmatrix} \Sigma & \sqrt{I - \Sigma^2} \\ \sqrt{I - \Sigma^2} & -\Sigma \end{pmatrix}
        \begin{pmatrix} V & 0 \\ 0 & I_n \end{pmatrix},
    \end{equation}
    它是 $A$ 的 $(1,1)$-block-encoding. 这说明 block-encoding 的模型本身不限制矩阵的结构,
    困难只在于 $U_A$ 是否能用简单电路实现.
\end{Eg}

\begin{Eg}[厄米矩阵的酉扩张]
    设 $A = A^\dagger$ 且 $\norm{A} \le 1$. 定义
    \begin{equation}
        U_A = \begin{pmatrix} A & \sqrt{I - A^2} \\ \sqrt{I - A^2} & -A \end{pmatrix}.
    \end{equation}
    由于 $A$ 是厄米的, $A$ 与 $\sqrt{I-A^2}$ 对易, 故 $U_A^\dagger U_A = I$, 即 $U_A$ 是
    酉矩阵, 且左上角块正是 $A$ —— 这是 $(1,1,0)$-block-encoding 的简单例子, 说明即使
    $A$ 本身非酉, 也可以作为更高维酉矩阵的子块出现.
\end{Eg}

\begin{Eg}[一般矩阵的酉扩张]
    对任意满足 $\norm{A} \le 1$ 的矩阵, 可构造
    \begin{equation}
        U_A = \begin{pmatrix} A & \sqrt{I - AA^\dagger} \\ \sqrt{I - A^\dagger A} & -A^\dagger \end{pmatrix},
    \end{equation}
    即任意 contraction($\norm{A} \le 1$)从纯数学角度都可以嵌入到更高维酉算子中
    (与上面的 SVD 构造等价). 然而量子算法真正重要的不是"这种酉扩张是否存在",
    而是能否用较少的量子门高效实现 $U_A$.
\end{Eg}

\begin{Def}[Hermitian block-encoding]
    若 $U_A$ 本身也是厄米矩阵, 则称其为 $A$ 的 $(\alpha,a,\varepsilon)$-Hermitian-block-encoding,
    记作 $U_A \in \mathrm{HBE}_{\alpha,a}(A,\varepsilon)$. Hermitian block-encoding
    为下一节 qubitization 提供了最简单的场景; 一般的 $s$-稀疏厄米矩阵可以用
    两个信号比特加 SWAP 的构造得到 $U_A \in \mathrm{HBE}_{s,n+2}(A)$(讲义命题 6.8).
\end{Def}

\begin{Note}[LCU 构造 block-encoding]
    线性组合酉(LCU, §3)是构造 block-encoding 的重要方法: 若 $A = \sum_j \alpha_j U_j$
    ($\alpha_j \ge 0$), 记 $\alpha = \sum_j \alpha_j$, 则取 $O_{prep}$ 与 $O_{sel}$ 后
    \begin{equation}
        U_A := (O_{prep}^\dagger \otimes I)\, O_{sel}\, (O_{prep} \otimes I)
    \end{equation}
    满足 $U_A \in \mathrm{BE}_{\alpha,a}(A)$. LCU 强调如何用 PREPARE 与 SELECT 构造
    酉算子的线性组合, 而 block-encoding 则把这一结构抽象为一般矩阵在高维酉矩阵中的表示方式.
\end{Note}

\subsection{基本运算}

Block-encoding 不仅是一种输入模型, 还构成一个可组合的\textbf{矩阵运算微积分}:
归一化因子 $\alpha$ 按照共轭转置、乘积、张量积与加法分别变化为
$\alpha$, $\alpha\beta$, $\alpha\beta$ 与 $\alpha + \beta$.

\begin{Prop}[Block-encoding 的基本运算]
    设 $U_A \in \mathrm{BE}_{\alpha,a}(A)$, $U_B \in \mathrm{BE}_{\beta,b}(B)$, 则
    \begin{enumerate}
        \item \textbf{共轭转置}: $U_A^\dagger \in \mathrm{BE}_{\alpha,a}(A^\dagger)$;
        \item \textbf{乘积}(同一系统寄存器上先后作用): 若两个 block-encoding 使用独立的
        辅助寄存器, 则
        \begin{equation}
            (U_A \otimes I_b)(I_a \otimes U_B) \in \mathrm{BE}_{\alpha\beta,\,a+b}(AB);
        \end{equation}
        \item \textbf{张量积}(不同系统寄存器上分别作用):
        \begin{equation}
            U_A \otimes U_B \in \mathrm{BE}_{\alpha\beta,\,a+b}(A \otimes B);
        \end{equation}
        \item \textbf{加法}(LCU 特例): 设两者辅助比特数相同(可补齐), 取
        \begin{equation*}
            O_{prep}\ket{0} = \sqrt{\frac{\alpha}{\alpha+\beta}}\,\ket{0} + \sqrt{\frac{\beta}{\alpha+\beta}}\,\ket{1},
        \end{equation*}
        \begin{equation*}
            O_{sel} = \ket{0}\bra{0}\otimes U_A + \ket{1}\bra{1}\otimes U_B,
        \end{equation*}
        则
        \begin{equation}
            (O_{prep}^\dagger \otimes I)\, O_{sel}\, (O_{prep} \otimes I)
            \in \mathrm{BE}_{\alpha+\beta,\,a+1}(A + B).
        \end{equation}
    \end{enumerate}
    这些规则都可以由定义直接验证; \textbf{后选择(post-selection)}成功概率相应地随 $\alpha$ 变化,
    因此 $\alpha$ 是 block-encoding 线路代价的核心参数.
\end{Prop}

\begin{Note}[矩阵多项式]
    乘法规则使 $A^2, A^3, \dots$ 的 block-encoding 可逐次得到(归一化因子相应为
    $\alpha^2, \alpha^3, \dots$), 加上加法规则, block-encoding 支持
    $A + B$, $AB$, $A^2$, $A^3$, $\dots$ 等基本矩阵代数运算, 自然引出
    \textbf{矩阵多项式(matrix polynomial)}
    \begin{equation}
        P(A) = c_0 I + c_1 A + c_2 A^2 + \cdots + c_d A^d.
    \end{equation}
    QSP 与 QSVT 的核心思想之一, 正是更加高效地实现这种矩阵多项式变换.
\end{Note}

\subsection{与奇异值、qubitization 的联系}

设 $A$ 的奇异值分解为 $A = \sum_j \sigma_j \ket{u_j}\bra{v_j}$, 其中 $\sigma_j \ge 0$,
$\ket{v_j}$, $\ket{u_j}$ 分别为右、左奇异向量. 由于 block-encoding 编码的是 $A/\alpha$,
归一化奇异值为
\begin{equation}
    x_j = \frac{\sigma_j}{\alpha} \in [0, 1].
\end{equation}
在后续 qubitization/QSVT 的分析中, 每个奇异值 $x_j$ 对应一个二维不变子空间, 其中
block-encoding 的作用转化为依赖于 $x_j$ 的二维旋转(定义角度 $\theta_j$ 使
$\cos\theta_j = \sigma_j/\alpha$, 或按约定取 $\sin\theta_j = \sigma_j/\alpha$):
\begin{equation}
    \text{singular value } \frac{\sigma_j}{\alpha}
    \quad\longleftrightarrow\quad
    \text{two-dimensional rotation angle } \theta_j.
\end{equation}
这是从 block-encoding 进入 qubitization、QSP 与 QSVT 的关键数学思想.

若 $A = A^\dagger$ 厄米, 则可作特征值分解 $A = \sum_j \lambda_j \ket{u_j}\bra{u_j}$; 当
$\norm{A} \le \alpha$ 时 $\lambda_j/\alpha \in [-1,1]$. 如果设计多项式
$P(x) \approx f(\alpha x)$, 则理想情况下可构造 $P(A/\alpha) \approx f(A)$ ——
这是现代量子矩阵函数算法的核心.

\subsection{与 OAA 的联系}

block-encoding 与 OAA(第 2 节)共享同一个成功子空间投影 $\Pi = \ket{0^a}\bra{0^a}\otimes I$;
OAA 中对应的反射算子为 $R_\Pi = 2\Pi - I$. 三者的作用可以概括为:
\begin{equation}
    \begin{aligned}
        \text{LCU}: &\quad \text{construct a desired linear combination},\\
        \text{block-encoding}: &\quad \text{represent a general matrix inside a unitary},\\
        \text{OAA}: &\quad \text{amplify the desired projected component}.
    \end{aligned}
\end{equation}

\subsection{与后续算法的联系}

block-encoding 并不是最终算法, 而是把矩阵转换为量子可操作对象的统一接口, 其典型发展路线为
\begin{equation}
    \text{LCU} \;\to\; \text{block-encoding} \;\to\; \text{qubitization} \;\to\;
    \text{QSP} \;\to\; \text{QSVT}.
\end{equation}
具体到几类重要问题(应用章将详述):
\begin{itemize}
    \item \textbf{Hamiltonian 模拟}: 对厄米 $H$, 由 block-encoding 给出 $\lambda/\alpha$
    在二维子空间结构中的编码, 用 qubitization 与 QSP 构造多项式 $P(x) \approx e^{-i\alpha t x}$,
    从而实现 $P(H/\alpha) \approx e^{-iHt}$;
    \item \textbf{线性方程组 (QLSP)}: 用 QSVT 构造 $P(x) \approx 1/x$, 近似实现
    $P(A) \approx A^{-1}$, 制备与 $A^{-1}\ket{b}$ 成比例的量子态
    $\ket{x} = A^{-1}\ket{b}/\norm{A^{-1}\ket{b}}$;
    \item \textbf{偏微分方程 (PDE)}: 许多 PDE 离散后写成线性系统 $A\mathbf{u} = \mathbf{f}$
    (如 Poisson 方程 $-\nabla^2 u = f$ 经有限差分/有限元离散), 量子路线为
    \begin{equation}
        \text{PDE discretization} \;\to\; A \;\to\; \text{block-encoding of } A
        \;\to\; \text{QSVT} \;\to\; A^{-1} \;\to\; \ket{u}.
    \end{equation}
    关键问题是如何利用 PDE 离散矩阵的结构(稀疏结构、Pauli 分解、LCU 分解、结构化有限差分
    矩阵、矩阵元 oracle、张量积结构等)高效构造 $A$ 的 block-encoding.
\end{itemize}

\begin{Note}[真正的算法成本]
    从数学上可以证明许多矩阵都存在酉扩张, 但量子算法真正关心的是 $U_A$ 能否高效实现.
    分析 block-encoding 时需要同时考虑: 归一化因子 $\alpha$、辅助比特数 $a$、近似误差
    $\varepsilon$、实现 $U_A$ 的门复杂度、以及调用 $U_A$ 与 $U_A^\dagger$ 的查询复杂度.
    因此仅仅写出 $U_A = \begin{pmatrix} A/\alpha & \ast \\ \ast & \ast \end{pmatrix}$
    并不意味着已经获得有效算法 —— 真正重要的是存在一个
    \textbf{efficiently implementable block-encoding}.
\end{Note}

\section{Qubitization 与 Chebyshev 多项式: 从谱变量到旋转与多项式逼近}

在 block-encoding 中, 一个一般的厄米算子 $H$ 被嵌入一个更大的酉算子:
\begin{equation}
    \big(\bra{0^a}\otimes I\big) U_H \big(\ket{0^a}\otimes I\big) = \frac{H}{\alpha},
    \qquad \alpha \ge \norm{H}.
\end{equation}
block-encoding 解决了"如何把一般矩阵表示为量子可实现的酉结构"的问题. 然而仅拥有
block-encoding 还不够: 在 Hamiltonian 模拟、QSP 与 QSVT 中, 还需要进一步把矩阵的谱信息
转化为一个简单的二维旋转结构. \textbf{Qubitization} 的核心作用正是:
\begin{equation}
    \text{把矩阵的特征值信息转化为二维不变子空间中的旋转角}.
\end{equation}
在这一二维旋转结构中, \textbf{Chebyshev 多项式}会自然出现: 特别地
$T_n(\cos\theta) = \cos(n\theta)$ 使量子旋转的重复作用与 Chebyshev 多项式之间建立了直接联系.
本节的核心数学链条为
\begin{equation}
    \frac{H}{\alpha} \;\longrightarrow\; \text{Qubitization} \;\longrightarrow\;
    x = \cos\theta \;\longrightarrow\; \cos(n\theta) = T_n(x) \;\longrightarrow\;
    P(x) \approx f(x) \;\longrightarrow\; f(H).
\end{equation}

\subsection{从 block-encoding 到谱变量}

设 $H = H^\dagger$ 是厄米 Hamiltonian, 谱分解为
\begin{equation}
    H = \sum_j \lambda_j \ket{\lambda_j}\bra{\lambda_j}.
\end{equation}
假设存在 $\alpha$-归一化 block-encoding $U_H$, 满足
$(\bra{0^a}\otimes I)U_H(\ket{0^a}\otimes I) = H/\alpha$. 由于 $\alpha \ge \norm{H}$,
对所有特征值有 $-1 \le \lambda_j/\alpha \le 1$. 定义\textbf{归一化谱变量(normalized spectral variable)}
\begin{equation}
    x_j := \frac{\lambda_j}{\alpha} \in [-1, 1],
    \qquad
    \frac{H}{\alpha}\ket{\lambda_j} = x_j \ket{\lambda_j}.
\end{equation}
Qubitization 的目的就是把实数 $x_j$ 转化为某个二维量子子空间中的旋转角.

\subsection{成功子空间投影与反射}

定义 block-encoding 的成功子空间投影
\begin{equation}
    \Pi = \ket{0^a}\bra{0^a}\otimes I,
    \qquad
    R_\Pi := 2\Pi - I.
\end{equation}
对成功子空间中的态 $R_\Pi\ket{\psi_{\mathrm{good}}} = \ket{\psi_{\mathrm{good}}}$, 对正交补中的态
$R_\Pi\ket{\psi_{\mathrm{bad}}} = -\ket{\psi_{\mathrm{bad}}}$. 这个反射与 OAA(第 2 节)中出现的
反射具有相同的数学结构 —— OAA 与 Qubitization 都大量利用
\begin{equation}
    \text{projection} \;\to\; \text{reflection} \;\to\; \text{two-dimensional invariant subspace}
\end{equation}
这一基本思想.

\subsection{标量启发: 反射到旋转}

对 $-1 \le \lambda \le 1$, 记 $\lambda = \cos\theta$, 考虑旋转矩阵
\begin{equation}
    O(\lambda) = \begin{pmatrix} \lambda & -\sqrt{1-\lambda^2} \\ \sqrt{1-\lambda^2} & \lambda \end{pmatrix}
    = \begin{pmatrix} \cos\theta & -\sin\theta \\ \sin\theta & \cos\theta \end{pmatrix}.
\end{equation}
其 $k$ 次幂为
\begin{equation}
    O(\lambda)^k = \begin{pmatrix} \cos(k\theta) & -\sin(k\theta) \\ \sin(k\theta) & \cos(k\theta) \end{pmatrix}
    = \begin{pmatrix} T_k(\lambda) & -\sqrt{1-\lambda^2}\, U_{k-1}(\lambda) \\
        \sqrt{1-\lambda^2}\, U_{k-1}(\lambda) & T_k(\lambda) \end{pmatrix},
\end{equation}
其中 $T_k$, $U_{k-1}$ 分别为第一、二类 Chebyshev 多项式
($T_k(\cos\theta) = \cos(k\theta)$, $U_{k-1}(\cos\theta) = \sin(k\theta)/\sin\theta$),
其定义与逼近性质将在本节后面系统介绍.
\textbf{Qubitization 的核心思想}正是: 若能把矩阵 $A$ 的每个本征方向上的反射结构
变为这样的旋转, 就立即得到 $T_k(A)$ 的 block-encoding.

\subsection{标准二维模型: 从反射构造 walk 算子}

先考虑标准的 Hermitian signal-unitary 情形: $U_H^\dagger = U_H$ 且 $U_H^2 = I$
(适当构造下, 一般 Hermitian block-encoding 可以转化为这种形式). 对归一化特征值
$x = \lambda/\alpha \in [-1,1]$, 在相应的二维不变子空间中 $U_H$ 表示为
\begin{equation}
    U_H(x) = \begin{pmatrix} x & \sqrt{1-x^2} \\ \sqrt{1-x^2} & -x \end{pmatrix},
    \qquad U_H(x)^2 = I,
\end{equation}
即 $U_H(x)$ 本身是反射型酉算子. 同时成功子空间反射在该二维子空间中为
$R_\Pi = \mathrm{diag}(1,-1)$. 定义\textbf{qubitized walk 算子(walk operator)}
\begin{equation}
    W := R_\Pi U_H
    \qquad\Longrightarrow\qquad
    W(x) = \begin{pmatrix} x & \sqrt{1-x^2} \\ -\sqrt{1-x^2} & x \end{pmatrix}.
\end{equation}
由于 $x \in [-1,1]$, 定义 $x = \cos\theta$($\theta = \arccos x$), 则
$\sqrt{1-x^2} = \sin\theta$, 于是
\begin{equation}
    W(x) = \begin{pmatrix} \cos\theta & \sin\theta \\ -\sin\theta & \cos\theta \end{pmatrix}
\end{equation}
正是一个二维旋转矩阵. 因此 Qubitization 实现了
\begin{equation}
    \frac{\lambda}{\alpha} = \cos\theta
    \quad\Longleftrightarrow\quad
    \text{Hamiltonian eigenvalue } \frac{\lambda}{\alpha}
    \;\longleftrightarrow\;
    \text{quantum rotation angle } \theta.
\end{equation}
注意这里 $W = R_\Pi U_H$ 与后文 Hermitian 情形中的 $O = U_A Z$ 相差一个转置
(旋转方向相反), 两者 $(1,1)$ 元均为 $T_k$ —— 约定不同, 结论相同.

\subsection{Hermitian block-encoding 的 qubitization}

设 $U_A \in \mathrm{HBE}_{1,a}(A)$, $A$ 的\textbf{谱分解(spectral decomposition)}为 $A = \sum_i \lambda_i \ket{v_i}\bra{v_i}$.
对每个本征态 $\ket{v_i}$,
\begin{equation}
    U_A \ket{0^a}\ket{v_i} = \lambda_i \ket{0^a}\ket{v_i} + \sqrt{1-\lambda_i^2}\,\ket{\Phi_i},
\end{equation}
其中 $\ket{\Phi_i}$ 是与所有 $\ket{0^a}\ket{\cdot}$ 正交的归一化态.
利用 $U_A$ 的厄米性再作用一次, 可得 $\mathrm{span}\{\ket{0^a}\ket{v_i}, \ket{\Phi_i}\}$
是 $U_A$ 的二维不变子空间, 且
\begin{equation}
    [U_A]_{B_i} = \begin{pmatrix} \lambda_i & \sqrt{1-\lambda_i^2} \\ \sqrt{1-\lambda_i^2} & -\lambda_i \end{pmatrix},
\end{equation}
即 $U_A$ 在该子空间上是\textbf{反射}. 投影算子 $\Pi = \ket{0^a}\bra{0^a}\otimes I_n$ 的表示为
$\mathrm{diag}(1,0)$, 故 $Z := 2\Pi - I$ 的表示为 $\mathrm{diag}(1,-1)$,
而\textbf{qubitization 迭代算子(qubitization iteration operator)}
\begin{equation}
    O = U_A Z
\end{equation}
的表示为旋转矩阵. 于是
\begin{equation}
    O^k = (U_A Z)^k, \qquad
    [O^k]_{B_i} = \begin{pmatrix} T_k(\lambda_i) & -\sqrt{1-\lambda_i^2}\,U_{k-1}(\lambda_i) \\
        \sqrt{1-\lambda_i^2}\,U_{k-1}(\lambda_i) & T_k(\lambda_i) \end{pmatrix},
\end{equation}
即 $(U_A Z)^k$ 是 $T_k(A)$ 的 $(1, a)$-block-encoding:
\begin{equation}
    \big(\bra{0^a}\otimes I_n\big) (U_A Z)^k \big(\ket{0^a}\otimes I_n\big) = T_k(A).
\end{equation}

\begin{Note}[$Z$ 的实现与电路]
    $Z$ 可以由信号比特与 $a$ 比特受控 $Z$ 门实现: 当辅助寄存器为 $\ket{0^a}$ 时施加 $-1$
    相位, 否则不变. 因此每次 qubitization 迭代只需 1 次 $U_A$ 与 $O(a)$ 个受控门.
\end{Note}

\subsection{一般 block-encoding 的 qubitization}

当 $U_A \notin \mathrm{HBE}$ 时, 对奇异(或本征)方向需要交替使用 $U_A$ 与 $U_A^\dagger$.
记 $U_A\ket{0^a}\ket{v_i} = \lambda_i\ket{0^a}\ket{w_i} + \sqrt{1-\lambda_i^2}\ket{\Phi_i}$
与对应的 $\ket{\Phi'_i}$, 则 $U_A$ 把二维子空间 $H_i$ 映射到 $H'_i$,
且两个子空间上 $Z$ 的表示相同. 于是

\begin{Thm}[一般 block-encoding 的 qubitization]
    设 $U_A \in \mathrm{BE}_{1,a}(A)$(不要求厄米). 定义
    \begin{equation}
        O := U_A Z U_A^\dagger Z,
    \end{equation}
    则 $O^k = (U_A Z U_A^\dagger Z)^k$ 是 $T_{2k}(A)$ 的 $(1,a)$-block-encoding,
    而 $U_A Z\, O^k$ 是 $T_{2k+1}(A)$ 的 $(1,a)$-block-encoding.
    电路上只需把 $U_A$, $U_A^\dagger$ 与 $Z$ 交替作用; Hermitian block-encoding
    ($U_A^\dagger = U_A$) 自动退化为上一小节的结构.
\end{Thm}

\begin{Proof}
    在每个二维子空间上 $U_A Z$ 与 $U_A^\dagger Z$ 都是旋转角 $\theta_i$
    ($\cos\theta_i = \lambda_i$) 的旋转, 故 $O = U_A Z U_A^\dagger Z$ 是旋转 $2\theta_i$,
    对应 $T_{2k}$; 奇数情形在末尾多乘一次 $U_A Z$ 即旋转 $(2k+1)\theta_i$.
\end{Proof}

\begin{Note}
    qubitization 的名字即来源于此: 每个本征(奇异)向量被"量子化"到一个二维不变子空间,
    大矩阵 $U_A$ 被部分分块对角化为 $N$ 个 $2\times2$ 块(也可以从 Jordan 引理或
    CS 分解的角度理解). 这一结构是后续 QSP/QSVT 的代数基础.
\end{Note}

\subsection{本征相位与二维不变子空间}

二维旋转矩阵 $W(x) = \begin{pmatrix} \cos\theta & \sin\theta \\ -\sin\theta & \cos\theta \end{pmatrix}$
的本征值为
\begin{equation}
    e^{\pm i\theta}.
\end{equation}
由于 $\cos\theta = \lambda/\alpha$, Hamiltonian 的特征值 $\lambda$ 被编码为 qubitized walk 的
本征相位 $\pm\theta$:
\begin{equation}
    \lambda \;\longrightarrow\; \frac{\lambda}{\alpha} \;\longrightarrow\;
    \theta = \arccos\left(\frac{\lambda}{\alpha}\right) \;\longrightarrow\; e^{\pm i\theta}.
\end{equation}
Qubitization 因此可以理解为一种\textbf{eigenvalue-to-eigenphase transformation}. 虽然完整
Hamiltonian 作用在高维 Hilbert 空间, 但每个本征态 $\ket{\lambda_j}$ 对应一个二维不变子空间
$\mathcal{K}_j$, 整个高维问题分解为许多彼此独立的二维旋转问题:
\begin{equation}
    \mathcal{H} \;\longrightarrow\; \bigoplus_j \mathcal{K}_j,
    \qquad
    W(x_j) = \begin{pmatrix} x_j & \sqrt{1-x_j^2} \\ -\sqrt{1-x_j^2} & x_j \end{pmatrix}.
\end{equation}
这正是后续 QSP 能同时对所有特征值进行函数变换的原因.

\subsection{重复 walk 与 Chebyshev 多项式}

连续执行 $n$ 次 walk 得到
\begin{equation}
    W(x)^n = \begin{pmatrix} \cos(n\theta) & \sin(n\theta) \\ -\sin(n\theta) & \cos(n\theta) \end{pmatrix},
\end{equation}
自然产生 $\cos(n\theta)$, $\sin(n\theta)$. 由于 $x = \cos\theta$, 关键问题是
$\cos(n\theta)$ 能否直接表示为 $x$ 的多项式 —— 答案正是第一类 Chebyshev 多项式.
\begin{equation}
    T_n(x) := \cos(n\arccos x), \qquad T_n(\cos\theta) = \cos(n\theta).
\end{equation}
于是 $[W(x)^n]_{11} = T_n(x)$, 投影回成功子空间后, $n$ 次 walk 的有效作用对应 $T_n(H/\alpha)$:
\begin{equation}
    \big(\bra{0^a}\otimes I\big) W^n \big(\ket{0^a}\otimes I\big) = T_n(H/\alpha)
    \quad\Longleftrightarrow\quad
    W^n \;\longleftrightarrow\; T_n(H/\alpha).
\end{equation}
这说明 Chebyshev 多项式不是人为加入 qubitization 的数学工具, 而是二维量子旋转重复作用后
自然产生的多项式结构. 为描述 $W^n$ 的非对角元, 还需定义\textbf{第二类 Chebyshev 多项式} $U_n$:
\begin{equation}
    U_n(x) := \frac{\sin((n+1)\theta)}{\sin\theta}, \qquad x = \cos\theta,
\end{equation}
由此 $\sin(n\theta) = \sqrt{1-x^2}\,U_{n-1}(x)$, $W^n$ 的完整矩阵形式为
\begin{equation}
    W(x)^n = \begin{pmatrix} T_n(x) & \sqrt{1-x^2}\,U_{n-1}(x) \\
        -\sqrt{1-x^2}\,U_{n-1}(x) & T_n(x) \end{pmatrix}.
\end{equation}

\subsection{Chebyshev 多项式的性质}

\begin{Def}[Chebyshev 多项式]
    第一类 Chebyshev 多项式 $T_k(x)$ 由 $T_k(\cos\theta) = \cos(k\theta)$ 定义,
    第二类 $U_k(x)$ 由 $U_k(\cos\theta) = \sin((k+1)\theta)/\sin\theta$ 定义
    ($k = 0, 1, 2, \dots$). 前几项为
    \begin{equation}
        T_0 = 1, \quad T_1(x) = x, \quad T_2(x) = 2x^2-1, \quad T_3(x) = 4x^3-3x,
        \qquad
        U_0 = 1, \quad U_1(x) = 2x.
    \end{equation}
    由三角恒等式二者都是实多项式, 并满足同一\textbf{三项递推(three-term recurrence)}
    \begin{equation}
        f_{k+1}(x) = 2x\, f_k(x) - f_{k-1}(x),
        \qquad
        T_0 = U_0 = 1, \quad T_1(x) = x, \quad U_1(x) = 2x,
    \end{equation}
    从而 $\deg T_k = \deg U_k = k$, 最高次项系数分别为 $2^{k-1}$ 与 $2^k$.
\end{Def}

\begin{Note}[有界性、正交性与极值性]
    ① \textbf{有界性}: 在 $[-1,1]$ 上 $|T_n(x)| \le 1$(由 $T_n(x) = \cos(n\arccos x)$),
    且在 $x = \cos(j\pi/k)$($j = 0, \dots, k$)处交替取 $\pm 1$, 端点 $T_k(\pm 1) = (\pm 1)^k$;
    ② \textbf{正交性}: 在权重 $1/\sqrt{1-x^2}$ 下
    \begin{equation}
        \int_{-1}^{1} \frac{T_i(x)T_j(x)}{\sqrt{1-x^2}}\,dx = 0 \qquad (i \ne j),
        \qquad
        \int_{-1}^{1} \frac{T_n(x)^2}{\sqrt{1-x^2}}\,dx =
        \begin{cases} \pi, & n = 0, \\ \pi/2, & n \ge 1; \end{cases}
    \end{equation}
    ③ \textbf{极值性(minimax property)}: 在首一多项式中 $2^{1-k}T_k$ 是
    $[-1,1]$ 上 $\sup$ 范数最小的, 极小值为 $2^{1-k}$ —— 因此 Chebyshev 多项式是
    $[-1,1]$ 上最优多项式逼近的基础.
\end{Note}

\begin{Note}[为什么用 Chebyshev 基而非幂基]
    普通幂基 $1, x, x^2, \dots$ 在高阶逼近中数值条件较差; Chebyshev 基则具有:
    $[-1,1]$ 上有界($|T_n| \le 1$)、正交性、near-minimax 逼近性质、对解析函数快速(指数型)收敛,
    以及与二维量子旋转 $\cos(n\theta)$ 的天然对应关系. 因此 Chebyshev 基同时具有良好的
    \textbf{逼近结构(approximation-theoretic structure)}与\textbf{量子旋转结构(quantum rotational structure)}.
\end{Note}

\subsection{Chebyshev 展开与逼近}

\begin{Note}[Chebyshev 展开与截断]
    定义在 $[-1,1]$ 上的适当函数 $f$ 可展开为 Chebyshev 级数
    \begin{equation}
        f(x) = \frac{c_0}{2} + \sum_{n=1}^{\infty} c_n T_n(x),
        \qquad
        c_n = \frac{2}{\pi}\int_{0}^{\pi} f(\cos\theta)\,\cos(n\theta)\,d\theta.
    \end{equation}
    截断到 $d$ 阶得 $P_d(x) = \frac{c_0}{2} + \sum_{n=1}^{d} c_n T_n(x) \approx f(x)$.
    光滑函数系数超代数衰减, 截断误差指数小; 这解释了应用章 Jacobi--Anger 展开中
    Bessel 系数 $J_n(t) \approx (t/2)^n/n!$ 的快速衰减(它正是在 $[-1,1]$ 上把 $e^{itx}$
    展开为 Chebyshev 多项式).
\end{Note}

\begin{Note}[一般谱区间与截断误差]
    若厄米矩阵 $A$ 的谱位于 $[a,b]$, 通过线性变换
    \begin{equation}
        x = \frac{2\lambda - (a+b)}{b-a},
        \qquad
        \widetilde A = \frac{2A - (a+b)I}{b-a},
    \end{equation}
    把谱映射到 $[-1,1]$; block-encoding 的归一化 $H/\alpha$ 本质上承担了类似的谱缩放.
    对可解析延拓到包含 $[-1,1]$ 的\textbf{Bernstein 椭圆} $E_\rho$($\rho > 1$)且在其中满足
    $|f(z)| \le M$ 的函数, Chebyshev 截断误差满足指数型界
    \begin{equation}
        \norm{f - P_d}_\infty \lesssim \frac{2M}{\rho - 1}\,\rho^{-d}.
    \end{equation}
    因此对解析函数, 达到误差 $\varepsilon$ 所需度数通常只随 $\log(1/\varepsilon)$ 增长 ——
    这是多项式量子算法获得较好精度复杂度的关键.
\end{Note}

\begin{Note}[带奇点的函数: 区间外逼近]
    对在 $[-1,1]$ 内部有奇点的函数(如 $x^{-1}$, 用于 QLSP), 不能直接在区间上逼近.
    方法是在奇点 $x=0$ 处留出邻域 $[-1/\kappa, 1/\kappa]$, 用奇多项式 $p$ 在
    $[1/\kappa,1]$ 上逼近 $\kappa/x$, 所需度数 $O(\kappa\log(\kappa/\varepsilon))$.
    这类"区间外逼近"是 QSVT 处理 $A^{-1}$、符号函数等的基本技术.
\end{Note}

\begin{Note}[矩阵函数的 Chebyshev 逼近]
    若 $f(\alpha x)$ 在 $[-1,1]$ 上被多项式 $P_d(x)$ 逼近($P_d(x) \approx f(\alpha x)$),
    则由函数演算
    \begin{equation}
        P_d(H/\alpha) = \sum_j P_d(\lambda_j/\alpha)\ket{\lambda_j}\bra{\lambda_j}
        \qquad\Longrightarrow\qquad
        P_d(H/\alpha) \approx f(H).
    \end{equation}
    由于 $f(H/\alpha) \approx \frac{c_0}{2}I + \sum_{n=1}^{d} c_n T_n(H/\alpha)$,
    而 qubitization 已给出 $T_n(H/\alpha) \leftrightarrow W^n$, 适当组合不同次数的
    qubitized walk 即可实现一般矩阵函数 $f(H)$ —— 这是从标量函数逼近进入矩阵函数逼近的关键一步.
\end{Note}

\subsection{Hamiltonian 模拟的 Chebyshev 表示}

Hamiltonian 模拟的目标是 $e^{-iHt}$. 令 $x = \lambda/\alpha$, 目标标量函数为
$f(x) = e^{-i\alpha t x}$. 利用\textbf{Jacobi--Anger 展开}($J_n$ 为第一类
\textbf{Bessel 函数}, $\tau = \alpha t$):
\begin{equation}
    e^{-i\tau x} = J_0(\tau) + 2\sum_{n=1}^{\infty} (-i)^n J_n(\tau) T_n(x)
    \quad\Longrightarrow\quad
    e^{-iHt} = J_0(\alpha t) I + 2\sum_{n=1}^{\infty} (-i)^n J_n(\alpha t) T_n(H/\alpha).
\end{equation}
截断到 $d$ 阶, $e^{-iHt} \approx J_0(\alpha t)I + 2\sum_{n=1}^{d} (-i)^n J_n(\alpha t) T_n(H/\alpha)$ ——
Hamiltonian 模拟被直接理解为一个 Chebyshev 多项式逼近问题. 与 Lie--Trotter 方法相比:
Trotter 把时间区间划分为小时间步
($e^{-i(H_1+H_2)t} \approx (e^{-iH_1t/r}e^{-iH_2t/r})^r$), 属于
\textbf{算子分裂 + 时间步进(operator splitting + time stepping)};
而 Qubitization/Chebyshev 路线 $e^{-iHt} \approx P_d(H/\alpha)$ 属于
\textbf{谱变换 + 多项式逼近(spectral transformation + polynomial approximation)}.
(应用章将给出截断阶数与复杂度的完整分析.)

\subsection{从 qubitization 到 QSP: 相干组合}

Qubitization 自然产生 $T_n(H/\alpha)$, 但这并不意味着简单执行 $W, W^2, \dots, W^d$
就自动实现了多项式 $P_d(H/\alpha) = \sum_n c_n T_n(H/\alpha)$ —— 不同项需要以正确的系数
$c_n$ 进行\textbf{相干组合(coherent combination)}. 一种方法是再次使用 LCU, 但这往往需要
额外辅助寄存器与振幅放大. \textbf{量子信号处理(quantum signal processing, QSP)}提供了更直接
强大的方法:
在 qubitized walk 之间插入一系列可调相位,
\begin{equation}
    U_{\vec\phi}(x) = e^{i\phi_0 Z} W(x)\, e^{i\phi_1 Z} W(x) \cdots W(x)\, e^{i\phi_d Z},
\end{equation}
通过适当选择相位序列 $\vec\phi = (\phi_0, \dots, \phi_d)$ 使某个矩阵元实现目标多项式 $P(x)$.
因此\textbf{Qubitization + Chebyshev 结构本身还不是 QSP, 而是 QSP 的数学基础}. QSP 所回答的问题是:
给定目标多项式 $P(x)$, 如何把它编译成一组 quantum phases? 可实现的 $P$ 需满足: 次数 $\le d$、
与序列长度匹配的奇偶性、在 $[-1,1]$ 上有界($|P(x)| \le 1$)、以及由酉性产生的实/虚部相容条件.
实际中分成两个子问题: \textbf{approximation problem}(找满足 QSP 条件的 $P(x) \approx f(x)$)与
\textbf{phase synthesis problem}(计算对应相位序列).

\subsection{从 QSP 到 QSVT}

QSP 研究标量变量 $x \in [-1,1]$ 上的变换 $x \to P(x)$; block-encoding 与 qubitization 使矩阵的
每个谱值(奇异值)都对应一个二维 signal subspace, 因此可对所有这些二维子空间同时应用同一套
QSP 相位序列. 对一般矩阵 $A = \sum_j \sigma_j \ket{u_j}\bra{v_j}$, 最终得到
$\sigma_j \to P(\sigma_j)$ —— 这就是\textbf{奇异值变换(quantum singular value transformation, QSVT)}.
知识结构可以概括为:
\begin{equation}
    \begin{aligned}
        \text{block-encoding} &\to \text{matrix inside a unitary},\\
        \text{qubitization} &\to \text{spectrum as rotation angles},\\
        \text{Chebyshev} &\to \text{rotation powers as polynomials},\\
        \text{QSP} &\to \text{programmable polynomial transformation},\\
        \text{QSVT} &\to \text{singular-value transformation}.
    \end{aligned}
\end{equation}
不同目标函数对应不同的量子矩阵算法:
$P(x) \approx e^{-itx}$(Hamiltonian 模拟)、$P(x) \approx 1/x$(量子线性系统)、
$P(x) \approx \mathrm{sign}(x)$(谱滤波)、一般 $P(x) \approx f(x)$(矩阵函数变换).

\subsection{与 OAA、Lie--Trotter 的联系}

\textbf{OAA 与 Qubitization}: 二者都使用投影与反射, 都归结为二维子空间中的旋转. 但 OAA 中
$\ket{\psi} = \sin\theta\ket{\psi_{\mathrm{good}}} + \cos\theta\ket{\psi_{\mathrm{bad}}}$,
核心目标是\textbf{放大成功振幅(amplitude amplification)};
Qubitization 中 $\lambda/\alpha = \cos\theta$, 核心目标是
\textbf{把谱信息编码进旋转角(spectral encoding)}. 数学结构相似, 功能不同.

\textbf{Qubitization 与 Lie--Trotter}: 二者目标都是 $e^{-iHt}$. Lie--Trotter 从分解
$H = \sum_j H_j$ 出发, 用 $e^{-iHt} \approx (\prod_j e^{-iH_jt/r})^r$, 误差主要来自非对易项
$[H_i, H_j]$, 属于\textbf{时间域算子分裂(time-domain operator splitting)};
Qubitization/QSP 先构造 $H/\alpha$ 的 block-encoding, 把特征值映射成旋转角, 再设计多项式
$P(x) \approx e^{-i\alpha tx}$, 属于\textbf{谱域多项式变换(spectral-domain polynomial transformation)}.

\section{QSP 与 QET: 相位因子编码多项式}

\subsection{量子信号处理: 标量表示定理}

在标量情形, $x \in [-1,1]$ 的 block-encoding 可以取为反射矩阵
\begin{equation}
    U_A(x) = \begin{pmatrix} x & \sqrt{1-x^2} \\ \sqrt{1-x^2} & -x \end{pmatrix},
\end{equation}
它是关于 $x$ 轴的反射. 乘以 $Z = \mathrm{diag}(1,-1)$ 后得到旋转
$O(x) = U_A(x)Z$. \textbf{量子信号处理(quantum signal processing, QSP)}考虑如下的
相位调制序列(也称 W 约定):
\begin{equation}
    U_\Phi(x) := e^{i\phi_0 Z}\, O(x)\, e^{i\phi_1 Z}\, O(x) \cdots e^{i\phi_{d-1} Z}\, O(x)\, e^{i\phi_d Z},
\end{equation}
其中 $\Phi = (\phi_0, \dots, \phi_d) \in \mathbb{R}^{d+1}$ 是\textbf{相位因子(phase factors)}.

\begin{Thm}[QSP 表示定理]
    对任意 $\Phi \in \mathbb{R}^{d+1}$, 存在多项式 $P, Q \in \mathbb{C}[x]$ 使得
    \begin{equation}
        U_\Phi(x) = \begin{pmatrix}
            P(x) & -Q(x)\sqrt{1-x^2} \\
            \overline{Q(x)}\,\sqrt{1-x^2} & \overline{P(x)}
        \end{pmatrix},
    \end{equation}
    且满足: (1) $\deg P \le d$, $\deg Q \le d-1$;
    (2) $P$ 的奇偶性为 $d \bmod 2$, $Q$ 的奇偶性为 $(d-1) \bmod 2$;
    (3) $|P(x)|^2 + (1-x^2)|Q(x)|^2 = 1$, $x \in [-1,1]$.
    反之, 任何满足这三条条件的 $P, Q$ 都可以由某个相位序列 $\Phi$ 表示.
\end{Thm}

\begin{Proof}[表示的正方向(归纳)]
    $d = 0$ 时 $U_\Phi = e^{i\phi_0 Z}$ 给出 $P = e^{i\phi_0}$, $Q = 0$. 设
    $U_{(\phi_0,\dots,\phi_{d-1})}$ 形如上式, 再右乘 $O(x)e^{i\phi_d Z}$ 并左乘 $e^{i\phi_0 Z}$,
    逐项展开可知新 $P$ 形如 $e^{i\phi_0}(xP - (1-x^2)Q)$, 度数至多比原来增 1,
    且奇偶性与 $d$ 相反 —— 归纳完成. 条件 (3) 是酉性的直接推论
    (矩阵的行范数平方和为 1).
\end{Proof}

\begin{Note}[实多项式版本的充分条件]
    在实际应用中最常用的是\textbf{实多项式}版本: 若 $P_{\mathrm{Re}} \in \mathbb{R}[x]$ 满足
    (1) 奇偶性为 $d \bmod 2$; (2) $|P_{\mathrm{Re}}(x)| \le 1$($x \in [-1,1]$),
    则存在相位序列 $\Phi$ 使 $U_\Phi$ 的 $(1,1)$ 元素恰为 $P_{\mathrm{Re}}(x)$.
    相比复数版本, 条件 (2) 只需对任意多项式做归一化缩放即可满足.
    相位因子的数值求解已有成熟的软件包(如 QSPPACK).
\end{Note}

\subsection{量子本征值变换 (QET)}

把标量 $x$ 换成矩阵 $A$, 即可把 QSP"提升"为矩阵版本 —— 这正是 qubitization 的意义:
每个本征方向上的二维旋转是同步进行的, 相位调制序列作用在辅助比特上, 与矩阵方向无关.

\begin{Thm}[QET: Hermitian block-encoding]
    设 $U_A \in \mathrm{HBE}_{1,a}(A)$. 对任意相位序列 $\Phi \in \mathbb{R}^{d+1}$,
    \begin{equation}
        U := e^{i\phi_0 Z}\, U_A\, e^{i\phi_1 Z}\, U_A \cdots e^{i\phi_{d-1} Z}\, U_A\, e^{i\phi_d Z}
    \end{equation}
    满足 $U = \begin{pmatrix} P(A) & \ast \\ \ast & \ast \end{pmatrix}$,
    即 $U \in \mathrm{BE}_{1,a}(P(A))$, 其中 $P$ 满足 QSP 定理的三条条件.
    电路结构为相位旋转 $e^{i\phi_j Z}$ 与 $U_A$ 交替作用 $O(d)$ 次.
\end{Thm}

\begin{Proof}
    qubitization 在每个不变子空间 $H_i$ 上把 $U_A$ 化为反射矩阵
    $[\lambda_i, \sqrt{1-\lambda_i^2}; \sqrt{1-\lambda_i^2}, -\lambda_i]$,
    恰好是标量情形 $U_A(\lambda_i)$ 的形式(可相差一个全局相位), 因此标量 QSP 定理逐块成立.
\end{Proof}

\begin{Eg}[Chebyshev 多项式与相位因子]
    取所有相位为零, 则 $U = (U_A Z)^d$, 恰好回到 qubitization 给出的
    $T_d(A)$ 的 block-encoding. 更一般地, 由于 $U_A e^{i\phi_j Z}$ 中的 $Z$ 可以吸收进
    相位因子($e^{i\phi Z}$ 与 $Z$ 的恒等式), QET 电路无需单独实现 $Z$ 门.
\end{Eg}

\begin{Eg}[实部多项式与一般多项式]
    QET 直接给出的是复多项式 $P(A)$. 若需要实多项式 $P_{\mathrm{Re}} = \Re P$,
    用 Hadamard 门叠加 $P(A)$ 与 $\overline{P}(A)$ 两个电路
    ($\overline{P}$ 由相位取负 $-(\phi_0,\dots,\phi_d)$ 实现), 得到
    $U_{P_{\mathrm{Re}}} \in \mathrm{BE}_{1,a+1}(P_{\mathrm{Re}}(A))$;
    这是线性组合 $[P(A) + \overline{P}(A)]/2$ 的实现. 对无确定奇偶性的一般实多项式
    $f = f_{\mathrm{even}} + f_{\mathrm{odd}}$, 分别构造后用一个辅助比特线性组合,
    得到 $U_f \in \mathrm{BE}_{2,a+2}(f(A))$; 复多项式情形则需要更深的 LCU 组合.
\end{Eg}

\begin{Note}[QET 与 LCU 的对比]
    QET 用 $O(d)$ 次 $U_A$ 与单比特旋转实现任意(满足奇偶条件的)多项式,
    辅助比特只增加 $O(1)$ 个, 也不需要 prepare/select oracle ——
    这正是前面 LCU 一节指出的两个缺点.
\end{Note}

\section{QSVT: 一般矩阵的奇异值变换}

\subsection{广义矩阵函数}

对一般(非厄米)矩阵 $A = W\Sigma V^\dagger$(奇异值分解), 定义三种\textbf{广义矩阵函数(generalized matrix functions)}:

\begin{Def}[广义矩阵函数]
    对定义在奇异值上的函数 $f: \mathbb{R} \to \mathbb{C}$,
    \begin{equation}
        f(A) := W f(\Sigma) V^\dagger, \qquad
        f^\triangleleft(A) := V f(\Sigma) V^\dagger, \qquad
        f^\triangleright(A) := W f(\Sigma) W^\dagger,
    \end{equation}
    分别称为(对称、左、右)\textbf{奇异值变换(singular value transformation)}. 三者之间有关系
    \begin{equation}
        f^\triangleleft(A) = f\big(\sqrt{A^\dagger A}\big), \qquad
        f^\triangleright(A) = f\big(\sqrt{AA^\dagger}\big).
    \end{equation}
\end{Def}

\subsection{一般矩阵的 qubitization 与 QSVT 定理}

对 $A = W\Sigma V^\dagger$, 设 $U_A \in \mathrm{BE}_{1,a}(A)$.
对右奇异向量 $\ket{v_i}$ 与左奇异向量 $\ket{w_i}$
($A\ket{v_i} = \lambda_i\ket{w_i}$, $A^\dagger\ket{w_i} = \lambda_i\ket{v_i}$),
$U_A$ 把二维子空间 $H_i = \mathrm{span}\{\ket{0^a}\ket{v_i}, \ket{\Phi_i}\}$ 映射到
$H'_i = \mathrm{span}\{\ket{0^a}\ket{w_i}, \ket{\Phi'_i}\}$, 且矩阵表示为反射.
因此与 Hermitian 情形完全平行地:

\begin{Thm}[一般矩阵的 qubitization]
    设 $U_A \in \mathrm{BE}_{1,a}(A)$(不需要厄米). 交替作用 $U_A$, $U_A^\dagger$ 与 $Z$,
    \begin{equation}
        (U_A Z U_A^\dagger Z)^k \in \mathrm{BE}_{1,a}\big(T_{2k}(A)\big), \qquad
        U_A Z (U_A Z U_A^\dagger Z)^k \in \mathrm{BE}_{1,a}\big(T_{2k+1}(A)\big),
    \end{equation}
    其中 $T_k(A) = W T_k(\Sigma) V^\dagger$ 是奇异值版本的 Chebyshev 变换.
\end{Thm}

\begin{Thm}[QSVT 定理(实多项式版本)]
    设 $A$ 由 $U_A \in \mathrm{BE}_{1,a}(A)$ 给出. 对满足
    (1) $\deg P_{\mathrm{Re}} = d$ 且奇偶性为 $d \bmod 2$;
    (2) $|P_{\mathrm{Re}}(x)| \le 1$, $x \in [-1,1]$ 的实多项式 $P_{\mathrm{Re}}$,
    存在相位序列 $\Phi \in \mathbb{R}^{d+1}$ 使 QSP/QET 型电路 $U$
    ($U_A$, $U_A^\dagger$ 交替, 共 $O(d)$ 次, 另用 $O(a)$ 个受控门与 $O(d)$ 个单比特旋转)满足:
    \begin{equation}
        d \text{ 为奇数}: \quad U \in \mathrm{BE}_{1,a+1}\big(P_{\mathrm{Re}}(A)\big),
        \qquad
        d \text{ 为偶数}: \quad U \in \mathrm{BE}_{1,a+1}\big(P_{\mathrm{Re}}^\triangleleft(A)\big).
    \end{equation}
\end{Thm}

\begin{Note}[奇偶性的意义]
    奇多项式 $P$($d$ 奇)把左、右奇异向量"配对", 实现 $W P(\Sigma) V^\dagger$;
    偶多项式 $P$($d$ 偶)只作用在右奇异向量上, 实现 $V P(\Sigma) V^\dagger$.
    由于奇异值非负, 任何目标函数 $f$ 都可以先做\textbf{奇(偶)延拓(odd/even extension)}再用 QSVT 实现:
    例如 $f(A)$($f$ 任意)可用 $f$ 的奇延拓经 $P_{\mathrm{Re}}(A)$ 实现.
\end{Note}

\begin{Eg}[Hermitian 情形]
    若 $A = A^\dagger$, 本征分解与奇异值分解通过 $\ket{w_i} = \mathrm{sgn}(\lambda_i)\ket{v_i}$
    联系. 对\textbf{偶}多项式 $P$: $P^\triangleleft(A) = P(A)$; 对\textbf{奇}多项式 $P$:
    $P(A) = P(A)$(此处 $P(A) := \sum_i P(\lambda_i)\ket{v_i}\bra{v_i}$ 为特征值函数).
    因此在 Hermitian 矩阵上 QET 与 QSVT 实现同一个矩阵函数, 只是视角不同.
\end{Eg}

\subsection{矩阵膨胀: QSP(Ā) 与 QSVT(A) 的等价}

\begin{Def}[厄米膨胀]
    对一般矩阵 $A$, 定义其\textbf{厄米膨胀}
    \begin{equation}
        \bar{A} := \begin{pmatrix} 0 & A \\ A^\dagger & 0 \end{pmatrix},
    \end{equation}
    由 $U_{\bar{A}} = \ket{0}\bra{1}\otimes U_A + \ket{1}\bra{0}\otimes U_A^\dagger$
    可得到 $\bar{A}$ 的 $(1,m)$-block-encoding(需要受控的 $U_A$, $U_A^\dagger$).
    $\bar{A}$ 的特征向量为
    \begin{equation}
        \ket{z_i^\pm} = \frac{1}{\sqrt{2}}\big(\ket{0}\ket{v_i} \pm \ket{1}\ket{w_i}\big),
        \qquad
        \bar{A}\ket{z_i^\pm} = \pm \lambda_i \ket{z_i^\pm}.
    \end{equation}
\end{Def}

\begin{Prop}[QET(Ā) 蕴含 QSVT(A)]
    对任意 $f \in \mathbb{C}[x]$, 记其偶部与奇部
    \begin{equation}
        f_{\mathrm{even}}(x) = \frac{f(x) + f(-x)}{2}, \qquad
        f_{\mathrm{odd}}(x) = \frac{f(x) - f(-x)}{2},
    \end{equation}
    则
    \begin{equation}
        f(\bar{A}) = \begin{pmatrix}
            f_{\mathrm{even}}^\triangleleft(A) & f_{\mathrm{odd}}(A)^\dagger \\
            f_{\mathrm{odd}}(A) & f_{\mathrm{even}}^\triangleright(A)
        \end{pmatrix}.
    \end{equation}
    特别地, 对偶函数 $f$, $f(\bar{A})\ket{0}\ket{\psi} = \ket{0}f_{\mathrm{even}}^\triangleleft(A)\ket{\psi}$;
    对奇函数 $f$, $f(\bar{A})\ket{0}\ket{\psi} = \ket{1}f(A)\ket{\psi}$ ——
    信号比特的测量结果是\textbf{确定性}的, 这正是 QSVT 不需要后选择(在偶函数情形)的原因.
\end{Prop}

\begin{Note}
    也就是说, 对 $\bar{A}$ 施加 QET 自动实现了对 $A$ 的 QSVT:
    框架因此可以统一处理厄米与非厄米矩阵 —— 求解非厄米线性方程组、一般哈密顿量模拟、
    伪逆与广义逆等都可以直接用 QSVT 实现.
\end{Note}

\section{应用}

\subsection{Hamiltonian 模拟: Jacobi--Anger 展开}

\begin{Thm}[基于 QET/QSVT 的哈密顿量模拟]
    设 $H$ 由 $U_H \in \mathrm{BE}_{1,a}(H)$ 给出($\norm{H} \le 1$). 利用
    \textbf{Jacobi--Anger 展开(Jacobi--Anger expansion)}($J_\nu$ 为第一类\textbf{Bessel 函数(Bessel function)}):
    \begin{equation}
        \cos(tx) = J_0(t) + 2\sum_{k\ge 1}(-1)^k J_{2k}(t)\, T_{2k}(x), \qquad
        \sin(tx) = 2\sum_{k\ge 0}(-1)^k J_{2k+1}(t)\, T_{2k+1}(x),
    \end{equation}
    截断到 $d = 2r + 1$ 阶, 其中
    \begin{equation}
        r = O\!\Big(t + \frac{\log(1/\varepsilon)}{\log\big(e + \log(1/\varepsilon)/t\big)}\Big),
    \end{equation}
    即可使 $\cos(tx) + i\sin(tx) = e^{itx}$ 的误差不超过 $\varepsilon$.
    分别用 QET 实现 $\cos$ 与 $\sin$ 部分(奇偶性恰好匹配)再用一个辅助比特组合,
    得到 $e^{itH}$ 的 $(2, m+2)$-block-encoding, 电路深度
    \begin{equation}
        O\!\Big(t + \frac{\log(1/\varepsilon)}{\log\log(1/\varepsilon)}\Big),
    \end{equation}
    其中每个 $U_H$ 调用一次. 该复杂度匹配哈密顿量模拟的查询下界
    $\Omega(t + \log(1/\varepsilon)/\log\log(1/\varepsilon))$ —— 相比 Trotter 方法
    $L = O(t^{1+1/p}\varepsilon^{-1/p})$ 的\textbf{乘性}依赖, 这里 $t$ 与
    $\log(1/\varepsilon)$ 是\textbf{加性}关系, 这是本质改进.
    若只有 $(\alpha,m)$-block-encoding(而非 $\alpha = 1$), 上式中的 $t$ 换成 $\alpha t$;
    后选择概率约为 $1/4$ 量级, 可再用 oblivious 振幅放大补偿.
\end{Thm}

\begin{Note}[为什么能避免时间分段]
    Jacobi--Anger 展开直接在整段时间 $[0,t]$ 上近似 $e^{itx}$, Bessel 系数
    $J_n(t) \approx (t/2)^n/n!$ 在 $n \gtrsim t$ 后超指数衰减, 因此截断阶数
    与 $t$ 成线性而非乘积关系; Trotter 型方法每步误差 $O(\tau^{p+1})$ 则把 $t$
    与 $\varepsilon$ 耦合在一起.
\end{Note}

\subsection{求解线性方程组 (QLSP)}

设 $A$ 可逆且 $\norm{A} \le 1$, \textbf{条件数(condition number)} $\kappa = \norm{A}\norm{A^{-1}}$,
奇异值位于 $[1/\kappa, 1]$. 目标函数 $f(x) = x^{-1}$ 在 $x = 0$ 奇异, 不能直接在
$[-1,1]$ 上做\textbf{多项式逼近(polynomial approximation)}, 因此用\textbf{奇多项式} $p$ 在区间外逼近:
\begin{equation}
    \Big|p(x) - \frac{\kappa}{x}\Big| \le \varepsilon, \qquad x \in [-1, -1/\kappa] \cup [1/\kappa, 1],
\end{equation}
且 $|p(x)| \le 1$($x \in [-1,1]$, 通过整体缩放因子实现). 这样的多项式存在且
$\deg p = O(\kappa \log(\kappa/\varepsilon))$ 量级. 用 QSVT(奇多项式, $d$ 奇)得到
$p(A) \in \mathrm{BE}_{1,a+1}(p(A))$.

\begin{Thm}[\textbf{量子线性系统问题(QLSP)} 的\textbf{查询复杂度(query complexity)}]
    设右端向量 $\ket{b}$ 由 oracle $U_b$ 制备, 记
    $\zeta := \norm{A^{-1}\ket{b}}$. 施加 QSVT 电路并测量辅助比特,
    成功概率为 $\Theta(\zeta^2/\kappa^2)$, 成功时系统态与归一化解
    $\ket{x} \propto A^{-1}\ket{b}$ 的误差为 $O(\varepsilon \kappa/\zeta)$;
    用振幅放大把成功概率提升到 $\Omega(1)$ 后, 对 $U_A$ 及其逆的总查询数为
    \begin{equation}
        O\!\Big(\frac{\kappa^2}{\zeta}\,\log\frac{\kappa}{\varepsilon}\Big),
    \end{equation}
    若 $\ket{b}$ 与最小奇异值对应的左奇异向量有 $\Omega(1)$ 重叠
    (此时 $\zeta = \Omega(1/\kappa)$), 则降为 $O(\kappa\log(1/\varepsilon))$;
    对 $U_b$ 的查询数为 $O(\kappa/\zeta)$.
\end{Thm}

\begin{Note}[与 HHL 的关系]
    QSVT 求解线性方程组是 HHL(下一模块)的现代替代: 不需要 QPE 的位数寄存器,
    不需要假设 $A$ 厄米(无需膨胀), 精度依赖为 $\log(1/\varepsilon)$ 而非
    $1/\varepsilon$ 量级. 两者都以 $\kappa$ 的(拟)线性或二次代价为特征.
\end{Note}

\subsection{其它应用与统一视角}

\begin{Note}[基态制备与特征值滤波]
    本注介绍\textbf{基态制备(ground-state preparation)}与\textbf{特征值滤波(eigenvalue filtering)}.
    设\textbf{谱隙(spectral gap)} $\Delta = E_1 - E_0 > 0$, 其中 $E_0$、$E_1$ 为基态与第一激发态能量.
    用偶多项式近似 $[0,1]$ 上的符号函数(在 $E_0$, $E_1$ 附近留出 $\Delta/2$ 的过渡带),
    度数 $O(\Delta^{-1}\log(1/\varepsilon))$ 即可. 施加 QET 后, 以概率
    $p_0(1-\varepsilon)$($p_0$ 为初态与基态的重叠)得到基态的近似,
    配合振幅放大后总查询数为
    $O(\Delta^{-1}p_0^{-1/2}\log(1/\varepsilon))$ —— 这是\textbf{变分量子本征求解器(variational quantum eigensolver, VQE)}与
    \textbf{绝热演化(adiabatic evolution)}的现代替代.
\end{Note}

\begin{Note}[Grover 搜索与固定点振幅放大的重访]
    在 QSVT 视角下, Grover 迭代正是对奇异值 $1/\sqrt{N}$ 的奇多项式变换
    ($p(x) \approx 1$ 于 $x \in [1/\sqrt{N}, 1]$); 单比特 QSP 电路统一了
    反射序列的几何. 更重要的是, \textbf{固定点振幅放大(fixed-point amplitude amplification)}: 用单调逼近 1 的
    多项式变换旋转角, 可以在\textbf{未知}成功概率 $p$ 的情形下使最终成功概率
    任意接近 1 —— 解决了 Module2 中普通振幅放大"需要知道 $p$"的根本缺陷.
    至此, QFT/QPE、Grover/AA、Trotter 与 HHL 等原语全部统一在
    block-encoding + qubitization + QSP/QSVT 的框架之下.
\end{Note}

\begin{Note}[本模块的主要文献]
    本模块的框架与定理主要来自:
    \begin{itemize}
        \item A. Gilyén, Y. Su, G. H. Low, N. Wiebe, \emph{Quantum singular value
        transformation and beyond}, STOC 2019; arXiv:1806.01838;
        \item G. H. Low, I. L. Chuang, \emph{Hamiltonian simulation by qubitization},
        Quantum 3, 163 (2019); arXiv:1610.06546;
        \item D. W. Berry, A. M. Childs, R. Cleve, R. Kothari, R. D. Somma,
        \emph{Simulating Hamiltonian dynamics with a truncated Taylor series},
        PRL 114, 090502 (2015);
        \item Lin Lin, \emph{Lecture Notes on Quantum Algorithms for Scientific
        Computation}(2022)以及 Lin Lin \& Nathan Wiebe 的同名讲义(2026).
    \end{itemize}
    相位因子的数值求解可参考 QSPPACK: \texttt{https://github.com/qsppack/QSPPACK}.
\end{Note}

\end{document}